%% kapa: exam=1 round=32 year=2021 period=1 form=A subject=economics subject_ko=경제학원론 seq=1 number=41
%% source: Q-Net 감정평가사 기출문제 (cst00309) artlSeq=5208397 file=extracted/1008/5208397/A형/2021년도 제32회 감정평가사 1차 1교시 A형.pdf
%% conversion: kordoc 4.24.51 -> tools/gen1.py (auto); review=proofread-sonnet; flags=none
\begin{question}{41}
무차별곡선에 관한 설명으로 옳지 않은 것은?
\begin{choices}
  \item 무차별곡선은 동일한 효용 수준을 제공하는 상품묶음들의 궤적이다.
  \item 무차별곡선의 기울기는 한계대체율이며 두 재화의 교환비율이다.
  \item 무차별곡선이 원점에 대해 오목하면 한계대체율은 체감한다.
  \item 완전대체재 관계인 두 재화에 대한 무차별곡선은 직선의 형태이다.
  \item 모서리해를 제외하면 무차별곡선과 예산선이 접하는 점이 소비자의 최적점이다.
\end{choices}
\end{question}
\answer{3}
